\documentclass[12pt]{article}
\usepackage[a4paper,margin=1in]{geometry}
\usepackage[T1]{fontenc}
\usepackage{lmodern,microtype}
\usepackage{amsmath,amssymb,amsthm,mathtools}
\usepackage{booktabs,array,enumitem,xcolor}
\usepackage[numbers,sort&compress]{natbib}
\usepackage[colorlinks=true,linkcolor=blue!55!black,citecolor=blue!55!black,urlcolor=blue!55!black]{hyperref}
\linespread{1.07}

\newtheorem{theorem}{Theorem}[section]
\newtheorem{proposition}[theorem]{Proposition}
\newtheorem{corollary}[theorem]{Corollary}
\theoremstyle{remark}
\newtheorem{remark}[theorem]{Remark}
\newcommand{\cL}{\mathcal L}
\newcommand{\norm}[1]{\left\lVert#1\right\rVert}

\title{A Canonical Biorthogonal Spectral Packet\\
\large Balanced Coordinates on Source--Observation Riesz Channels}
\author{Liang Wang}
\date{July 2026}

\begin{document}
\maketitle

\begin{abstract}
RH-195 associates each selected simple physical eigenvalue with exact right
and left matrix states $X_i=P_iS$ and $Y_i=P_i^*O^*$ whose cross pairing is
diagonal with entries equal to transfer residues.  This paper constructs
optimal packet coordinates on the resulting right and left spectral spaces.

Let $E_R=\operatorname{span}\{X_i\}$ and
$E_L=\operatorname{span}\{Y_i\}$, assume every residue is nonzero, and let
$Q_R,Q_L$ be orthonormal bases.  If
$Q_L^*Q_R=U\Sigma V^*$, define
\[
 V_c=Q_RV\Sigma^{-1/2},
 \qquad
 W_c=Q_LU\Sigma^{-1/2}.
\]
Then $W_c^*V_c=I$, both frame norms equal
$\sigma_{\min}^{-1/2}$, and their norm product
$\sigma_{\min}^{-1}$ is optimal among all biorthogonal coordinates on the
same two spaces.

Because $E_R$ and $E_L$ are exact invariant spaces,
\[
 \mathcal L_AV_c=V_cK_c,
 \qquad
 \mathcal L_A^*W_c=W_cK_c^*.
\]
Both directed residuals vanish, $K_c$ is similar to the selected physical
spectral restriction, and its determinant and every power trace are exact.
A 140-case complex nonnormal audit verifies all identities with zero
failures.

The packet is canonical relative to the chosen Riesz contours, source, and
observation.  It removes the complement self-energy inside that selected
channel.  Intrinsic contour selection, interval validation, and uniform
cross-level transversality remain open.
\end{abstract}

\section{From residue coordinates to balanced coordinates}

For pairwise disjoint simple eigenvalues, RH-195 gives right and left channel
states
\begin{equation}\label{eq:channel-states}
 X_i=P_iS,
 \qquad Y_i=P_i^*O^*,
\end{equation}
with
\begin{equation}\label{eq:residue-cross}
 \langle Y_i,X_j\rangle_F=\delta_{ij}r_i.
\end{equation}
If $r_i\ne0$, the direct normalization
$V_i=X_i$, $W_i=Y_i/\overline{r_i}$ is biorthogonal.  It preserves mode
labels and diagonalizes the compressed operator, but it need not balance
coordinate norms \citep{WangRH195}.  The cross-Gram balancing used below is
the exact spectral specialization of the temporal construction in
RH-184 \citep{WangRH184}.

The geometry depends only on the two subspaces
\begin{equation}\label{eq:spaces}
 E_R=\operatorname{span}\{X_1,\ldots,X_r\},
 \qquad
 E_L=\operatorname{span}\{Y_1,\ldots,Y_r\}.
\end{equation}
Nonzero residues imply that the pairing between them is nondegenerate.

\section{Cross-Gram geometry}

Choose Frobenius-orthonormal bases
$Q_R,Q_L\in\mathbb C^{N\times r}$ for the vectorized spaces and form
\begin{equation}\label{eq:cross-gram}
 H=Q_L^*Q_R.
\end{equation}
The singular values of $H$ are the cosines of the principal angles between
$E_L$ and $E_R$.  Transversality is equivalent to
$\sigma_{\min}(H)>0$.

Let
\begin{equation}\label{eq:svd}
 H=U\Sigma V^*.
\end{equation}

\begin{theorem}[Balanced spectral packet]\label{thm:balanced}
Define
\begin{equation}\label{eq:balanced-frames}
 V_c=Q_RV\Sigma^{-1/2},
 \qquad
 W_c=Q_LU\Sigma^{-1/2}.
\end{equation}
Then
\begin{equation}\label{eq:biorthogonal}
 W_c^*V_c=I,
\end{equation}
and
\begin{equation}\label{eq:frame-norms}
 \norm{V_c}=\norm{W_c}=\sigma_{\min}(H)^{-1/2}.
\end{equation}
\end{theorem}

\begin{proof}
Insert the singular-value decomposition into
$W_c^*V_c$.  Since $Q_R,Q_L,U,V$ are isometries, the operator norms of the
frames equal the largest diagonal entry of $\Sigma^{-1/2}$.
\end{proof}

\section{Optimality}

Any frames $V,W$ spanning $E_R,E_L$ and satisfying $W^*V=I$ define the same
oblique projector $\Pi=VW^*$ from the ambient space onto $E_R$ along
$E_L^\perp$.  Its norm is
\begin{equation}\label{eq:projector-norm}
 \norm\Pi=\sigma_{\min}(H)^{-1}.
\end{equation}

\begin{theorem}[Optimal norm product]\label{thm:optimal}
Every biorthogonal realization on $(E_R,E_L)$ obeys
\begin{equation}\label{eq:product-lower}
 \norm V\norm W\ge\sigma_{\min}(H)^{-1}.
\end{equation}
The balanced frames attain equality.  They also minimize
$\max(\norm V,\norm W)$, whose optimal value is
$\sigma_{\min}(H)^{-1/2}$.
\end{theorem}

\begin{proof}
Submultiplicativity gives
$\norm{VW^*}\le\norm V\norm W$, while the left side equals
$\sigma_{\min}^{-1}$.  The balanced norms attain the lower bound.  For any
positive numbers with product at least $c$, their maximum is at least
$\sqrt c$; scalar gauge balancing attains this value.
\end{proof}

The potentially large conditioning is intrinsic to the source-observation
subspaces.  No alternative basis can remove it.

\section{Exact two-sided invariance}

The right space is invariant under $\cL_A:X\mapsto AX$, and the left space
is invariant under $\cL_A^*$.  Hence there are matrices $K_R,K_L$ with
\begin{equation}\label{eq:invariance}
 \cL_AV_c=V_cK_R,
 \qquad
 \cL_A^*W_c=W_cK_L.
\end{equation}

\begin{theorem}[One exact compressed operator]\label{thm:exact-compression}
With
\begin{equation}\label{eq:compressed}
 K_c=W_c^*\cL_AV_c,
\end{equation}
one has
\begin{equation}\label{eq:two-sided}
 \boxed{
 \cL_AV_c=V_cK_c,
 \qquad
 \cL_A^*W_c=W_cK_c^*.
 }
\end{equation}
Thus both directed Ritz residuals vanish exactly.
\end{theorem}

\begin{proof}
The first invariant relation and $W_c^*V_c=I$ imply $K_R=K_c$.  Taking the
adjoint of the left invariant relation gives
$W_c^*\cL_A=K_L^*W_c^*$.  Multiplying by $V_c$ shows
$K_L^*=K_c$, hence the second identity.
\end{proof}

This is the exact endpoint that the temporal bi-Krylov packet approximates.
In these coordinates the packet-to-complement couplings are zero because
the selected subspaces are reducing in the biorthogonal sense.

\section{Spectrum, determinant, and traces}

In the residue-labeled frame, the compressed matrix is
$\Lambda=\operatorname{diag}(\lambda_1,\ldots,\lambda_r)$.  Every other
biorthogonal frame on the same spaces differs by a gauge
$V\mapsto VG$, $W\mapsto W(G^{-1})^*$, so
\begin{equation}\label{eq:gauge-similarity}
 K\mapsto G^{-1}KG.
\end{equation}

\begin{corollary}[Exact finite spectral ledger]\label{cor:ledger}
For the balanced packet,
\begin{align}
 \det(zI-K_c)&=\prod_{j=1}^r(z-\lambda_j),\label{eq:determinant}\\
 \operatorname{tr}K_c^q&=\sum_{j=1}^r\lambda_j^q
 \quad(q\ge1).\label{eq:traces}
\end{align}
\end{corollary}

These are exact finite channel identities.  They are not the determinant or
trace of the complete Frobenius operator, whose multiplicities remain
multiplied by $m$.

\section{Canonicity and its limits}

The spaces $E_R,E_L$ are determined by $(A,S,O)$ and the selected Riesz
contours.  The balanced packet is unique up to unitary transformations within
repeated singular-value blocks and harmless phase conventions.  Its
determinant and traces are fully invariant under this residual gauge.

This is a useful but conditional canonicity.  If the contours were selected
by first inspecting the target eigenvalues, the construction has not yet
explained why those modes are dynamically privileged.  Gate A needs an
a priori selection principle, for example an edge-gap, source-history, or
renormalization rule that survives refinement.

\section{Identity audit}

The numerical audit uses 140 complex nonnormal systems with base dimensions
4--10 and source widths 1--4.  Four spectral channels of largest modulus are
formed from exact computed left/right eigendata.  The code balances their
right and left matrix-state spaces and verifies:
\begin{enumerate}[leftmargin=2em]
\item equality of balanced frame norms;
\item biorthogonality;
\item zero right and left invariant residuals;
\item equality of compressed and selected spectra;
\item Newton power-trace identities through power eight.
\end{enumerate}
All cases pass the declared $10^{-8}$ tolerance.

\section{Position in the route}

The local route now contains an exact endpoint:
\begin{equation*}
 \text{temporal packet}
 \longrightarrow
 \text{matched Riesz channels}
 \longrightarrow
 \text{balanced exact packet}.
\end{equation*}
The first arrow is finite and numerical; the second is exact once the
contours are fixed.  RH-197 measures the physical residues and the intrinsic
cross-angle condition of the matched quartet.  RH-198 then quantifies how
the temporal packet approaches this exact endpoint.

No statement here constructs a self-adjoint operator, a Hilbert--P\'olya
spectrum, a prime-power trace formula, or a zeta divisor.  The work remains
inside the physical spectral interface portion of Gate A.

\bibliographystyle{plainnat}
\bibliography{references}
\end{document}
